%── themes/story-accents.tex ────────────────────────────────────────────────
%
% Paleta tonal del poster: grises con tinte azul-índigo.
% Son los overlays de Catppuccin Mocha usados sobre blanco —
% más sofisticados que grises puros, coherentes con el sistema Ctp.
%
% Sistema de acento: un único color por historia, definido en
% story-content.tex como \StoryAccentColor (nombre de color Ctp*).
% Se materializa en \begin{document} con:
%   \colorlet{StoryAccent}{\StoryAccentColor}
%   \colorlet{StoryAccentFaint}{\StoryAccentColor!12!white}
%───────────────────────────────────────────────────────────────────────────

% ─── Escala tonal (4 niveles + fantasma) ────────────────────────────────────
\definecolor{InkBlack}{HTML}{1e1e2e}   % héroe, énfasis máximo
\definecolor{InkDark} {HTML}{313244}   % títulos principales
\definecolor{InkMid}  {HTML}{585b70}   % body text, subtítulos
\definecolor{InkLight}{HTML}{9399b2}   % metadatos, captions, fechas
\definecolor{InkFaint}{HTML}{cdd6f4}   % líneas sutiles, número decorativo

% ─── Nota: NO se sobreescriben colores Ctp* aquí ────────────────────────────
% La paleta Catppuccin completa sigue disponible para ecuaciones
% (equation-styles.tex) y para el acento de color de la historia.